\documentclass[10pt,a4paper]{amsart}
\usepackage[margin=19mm]{geometry}
\usepackage{amsmath,amssymb,mathtools}
\usepackage[scr=rsfs,cal=euler]{mathalfa}
\usepackage{xfrac}
\usepackage[unicode,hidelinks,bookmarks=false]{hyperref}
\newcommand{\ie}{i.e.\ }
\newcommand{\cf}{cf.\ }
\newcommand{\resp}{respectively\ }
\newcommand{\Subsection}{\subsection}
\providecommand{\nobreakdash}{\nobreak\hbox{-}}
\setlength{\emergencystretch}{2em}
\multlinegap0pt
\usepackage[shortlabels]{enumitem}
\usepackage{braket}
\usepackage{amssymb}
\usepackage{dsfont}
\usepackage{tikz}
\usepackage{xcolor}
\newtheorem{theorem}{Theorem}
\title{Axiomatic Foundation of Quantum-Inspired Distance Metrics}
\author{Maryam Bagherian}
\date{Source: arXiv:2603.00406v1 (CC BY 4.0)}
\begin{document}
\maketitle

\noindent\emph{Source and licence.} Axiomatic Foundation of Quantum-Inspired Distance Metrics. Version arXiv:2603.00406v1 (CC BY 4.0), distributed by arXiv under the Creative Commons Attribution 4.0 International licence, http://creativecommons.org/licenses/by/4.0/, as recorded in the arXiv licence metadata for this exact version and shown on the abstract page https://arxiv.org/abs/2603.00406v1. Full licence notice: \texttt{licenses/arxiv-2603.00406-CC-BY-4.0.txt}.\par\smallskip

\noindent\textit{Editorial scope.} Counted unit: the article's theorem thm:entanglement-bounds (source0.tex lines 1657--1669). The article's complete original proof of it is delivered with it (source0.tex lines 1671--1719, one \texttt{proof} environment of 49 lines). No mathematics is written, repaired or re-derived: the statement and the proof are the article's own lines, in the article's own notation, and no retained line is substituted or re-flowed.  No further source-local context is needed: every cross-reference of every retained line resolves inside the two delivered spans.  Nothing was omitted from any retained span, so the package carries no omission record.  No retained line contains a citation, so the wrapper carries no bibliography and no bibliography entry.  No retained line contains a figure environment, an image-inclusion command or drawing code, so no image is delivered and the package declares no asset dependency.  Editorial formatting only, in addition: an AMS \texttt{amsart} preamble that restates the article's own declarations and macros used by the retained lines, namely the environments \texttt{theorem} on separate counters, together with the standard packages those lines need.  The retained environments print in this document as Theorem 1 in order of appearance, whereas the article prints the same items with the numbering of its own full document.  The complete native source of the article as distributed in the arXiv e-print for this exact version is bundled unchanged as \texttt{sources/195490/source0.tex}.
\begin{theorem}[Entanglement Distance Bounds]
\label{thm:entanglement-bounds}
Let 
\[
d_E(\ket{\psi},\ket{\varphi}) = \sqrt{d_{\text{FS}}^2(\ket{\psi},\ket{\varphi}) + |E(\ket{\psi}) - E(\ket{\varphi})|^2},
\] 
where $E$ denotes the entanglement entropy. Then:
\begin{enumerate}[label=(\roman*)]
    \item $d_{\text{FS}} \le d_E \le d_{\text{FS}} + |E(\ket{\psi}) - E(\ket{\varphi})|$,
    \item $d_E$ satisfies the triangle inequality,
    \item $d_E$ is entanglement-aware: states with identical reduced density matrices can have $d_E>0$.
\end{enumerate}
\end{theorem}

\begin{proof}
{(i) Bounds:} Apply the Pythagorean inequality
\[
a \le \sqrt{a^2 + b^2} \le a + b, \quad \forall a,b \ge 0,
\]
with $a = d_{\text{FS}}(\ket{\psi},\ket{\varphi})$ and $b = |E(\ket{\psi}) - E(\ket{\varphi})|$.\\
{(ii) Triangle inequality:}
Define
\[
\mathbf{v}_{\psi\varphi}
=
\big(
d_{\mathrm{FS}}(\ket{\psi},\ket{\varphi}),
\, |E(\ket{\psi}) - E(\ket{\varphi})|
\big)
\in \mathbb{R}^2.
\]
Then
\[
d_E(\ket{\psi},\ket{\varphi})
=
\|\mathbf{v}_{\psi\varphi}\|_2.
\]
Using the triangle inequality for $d_{\mathrm{FS}}$ and for absolute values,
\[
\mathbf{v}_{\psi\chi}
\le
\mathbf{v}_{\psi\varphi}
+
\mathbf{v}_{\varphi\chi},
\]
componentwise. The Minkowski inequality in $\mathbb{R}^2$ then implies
\[
d_E(\ket{\psi},\ket{\chi})
\le
d_E(\ket{\psi},\ket{\varphi})
+
d_E(\ket{\varphi},\ket{\chi}).
\]
{(iii) Entanglement awareness:} Consider the Bell states
\[
\ket{\Phi^\pm} = \frac{1}{\sqrt{2}}(\ket{00} \pm \ket{11}).
\]
Both have identical reduced density matrices $\rho_A = \rho_B = \frac{1}{2} \mathbb{I}_2$, yet
\begin{multline}\notag 
d_E(\ket{\Phi^+},\ket{\Phi^-}) \\= \sqrt{(\pi/2)^2 + |E(\ket{\Phi^+}) - E(\ket{\Phi^-})|^2} = \pi/2 > 0,
\end{multline}
so $d_E$ distinguishes global states even when local reductions are identical.
\end{proof}

\end{document}
